	In this section, we explain our model of a self-modifying agent, which is borrowed from \citet{everitt}. We will extend this model to include bounded rationality in \Cref{sec:def_bounded_rat}.

	We use a modified version of the \textit{standard cybernetic model}. In this model, an agent interacts with the environment in discrete time steps. At each time step $t$, the agent performs an \textit{action} $a_t$ from a finite set $\mathcal{A}$ and the environment responds with a \textit{perception} $e_t$ from a finite set $\mathcal{E}$. An \textit{action-perception pair} $\mae_t$ is an action concatenated with a perception. A \textit{history} is a sequence of action-perception pairs $\mae_1\mae_2...\mae_t$. We will often abbreviate such sequences to $\mae_{< t}=\mae_1...\mae_{t-1}$ or $\mae_{n:m}=\mae_n...\mae_m$. A \textit{complete} history $\mae_{1:\infty}$ is a history containing information about all the time steps.
%	
%	\[diagram\] \jakub{za mě není nutnost. Neměl jsem pocit, že by mi ten diagram v Everittovo článku nějak pomohl s pochopením. (Tedy, za předpokladu, že to je ten diagram, co bys tu chtěl)}
%	

    An agent can be described by its policy $\pi$. The policy\footnote{For a set $S$, $S^*$ denotes the set of finite sequences of elements from $S$} $\pi \colon (\mathcal{A \times E})^* \to \mathcal{A}$ is used to determine the agent's next action from the history at time $t$. We consider (bounded-rational) utility maximizers, where the policy is (partially) determined by the instantaneous utility function $u$, belief $\rho$ and discount factor $\gamma$. We sometimes use the notation $\kappa = (u, \rho, \gamma)$, where $\kappa$ is called the agent's \emph{knowledge}. The utility function $\tilde{u} \colon (\mathcal{A \times E})^\infty \to \mathbb{R}$ describes how much the agent prefers the complete history $\mae_{1:\infty}$ compared to other complete histories. We will assume that the total utility is a discounted sum of instantaneous utilities given by the instantaneous utility function $u \colon (\mathcal{A \times E})^* \to [0,1]$. Formally, $ \tilde{u}(\mae_{1:\infty}) = \sum ^{\infty }_{t=1}\gamma^{t-1} u(\mae_{\leq t})$. The discount factor $\gamma$ describes how much the agent prefers immediate reward compared to the same reward at a later time. Smaller $\gamma$ means heavier discounting of the future and stronger preference for immediate reward. Note that the maximum achievable utility is $\frac{1}{1-\gamma}$, which happens when $u(\mae_t)=1$ at each step. Also note that instantaneous utility depends not only on the latest perception but can also depend on all previous perceptions and actions.
	
	In addition to all this, an agent has a belief $\rho \colon (\mathcal{A \times E})^* \times \mathcal{A} \to \Delta \bar{\mathcal{E}}$ where $\Delta \bar{\mathcal{E}}$ is the set of full-support probability distributions over $\mathcal{E}$. This is a function which maps any history ending with an action onto a probability distribution over the next perceptions. Intuitively speaking, the belief describes what the agent expects to see after it performs an action given a certain history. A belief together with a policy induce a measure on the set $(\mathcal{A \times E})^*$ using $P(e_t \mid \mae_{< t}a_t) = \rho(e_t \mid \mae_{< t}a_t)$  and $P(a_t \mid \mae_{< t}) = 1$ if $\pi(\mae_{<t})=a_t$ and 0 otherwise. Intuitively speaking, this probability measure captures probabilities assigned by the agent to possible futures.
	
	Following the reinforcement learning literature, we define the \textit{value function} $V\colon (\mathcal{A \times E})^* \to \mathbb{R}$ as the expected future discounted utility:
	\[ V^\pi(\mae_{<t})=\mathbb{E}[\ \sum_{t'=t}^\infty \gamma^{t'-t}u(\mae_{<t'}) \ ]\]
	The expectation value on the right is calculated with respect to belief $\rho$ and assuming the agent will follow the policy $\pi$. Intuitively, the value function describes how promising the future seems. When the value \(V^\pi(\mae_{<t})\) of a history is high, it means we can expect an agent with policy $\pi$ to collect a lot of utility in the future starting from this history. Note that when calculating V-values, instantaneous utilities are multiplied by \(\gamma^{t'-t}\)  rather than  \(\gamma^{t'}\). This means that V-values can remain high throughout the whole history and are not affected by discounting. We define the Q-value of an action as the expected future discounted utility after taking that action:
	\[ Q^\pi(\mae_{<t} a_t)=\mathbb{E}[u(\mae_{1:t}) + V^{\pi}(\mae_{1:t})] \]
	where the expectation is over the next perception drawn from the belief (note that belief is a probability distribution)

	The Q-value measures how good an action is given that the agent will later follow policy $\pi$. A policy $\pi^*$ is an \textit{optimal policy} when $V^{\pi^*}(\mae_{<t}) = \sup_{\pi} V^{\pi}(\mae_{<t})$ for all histories $\mae_{<t}$ (such a policy always exists, as shown in \cite{Lattimore2014}). 	
	\subsection{Self-modification model}
	In this section, we extend the formalism above to include the possibility of self-modification. Since we are interested in the worst-case scenario, we assume the agent has unlimited self-modification ability. Worst-case results derived for such an agent will also hold for agents with limited ability to self-modify.
	
	\begin{definition}
		A policy self-modification model is defined as a quadruple $(\breve{\mathcal{A}}, \mathcal{E, P}, \iota)$ where $\breve{\mathcal{A}}$ is the set of \textit{world actions}, $\mathcal{E}$ is the set of \textit{perceptions}, $\mathcal{P}$ is a non-empty set of \textit{names} and $\iota$ is a map from $\mathcal{P}$  to the set of all policies $\Pi$. 
	\end{definition}
	At every time step, the agent chooses an action $a_t = (\breve{a}_t, p_{t+1}
	)$ from the set $\mathcal{A} = (\breve{\mathcal{A}} \times \mathcal{P})$. The first part $\breve{a}_t$ describes what the agent ``actually does in the world'' while the second part chooses the policy $\pi_{t+1}=\iota(p_{t+1})$ for the next time step. We will also use the notation  $a_t = (\breve{a}_t, \pi_{t+1}
	)$, keeping in mind that only policies with names may be chosen. This new policy is used in the next step to pick the action $a_{t+1}=\pi_{t+1}(\mae_{1:t})$. Note that only policies with names can be chosen and that $\mathcal{P} = \Pi$ is not a possibility because it entails a contradiction: 
	$ |\Pi| = |(\breve{\mathcal{A}} \times \mathcal{E}\times \Pi)|^{|(\breve{\mathcal{A}} \times \mathcal{E}\times \Pi)^*|} > 2^{|\Pi|}>|\Pi|$. A history can now be written as:
	\[\mae_{1:t} = a_1 e_1 a_2 e_2 ... a_t e_t = \breve{a}_1 \pi_2 e_1 \breve{a}_2 \pi_3 e_2 ... \breve{a}_t \pi_{t+1} e_t \]
	The subscripts for policies are one time step ahead because the policy chosen at time $t$ is used to pick an action at time $t+1$. The subscript denotes at which time step the policy is used. Policy $\pi_t$ is used to choose the action $a_t = (\breve{a_t}, \pi_{t+1})$. No policy modification happens when $a_t = (\breve{a_t}, \pi_t)$.
	%\[diagram \  2\]

	In the previous section, we used these rules to calculate the probability of any finite history: $P(e_t \mid \mae_{< t}a_t) = \rho(e_t \mid \mae_{< t}a_t)$  and $P(a_t \mid \mae_{< t}) = 1$ if $\pi(\mae_{<t})=a_t$ and 0 otherwise. However, the second rule doesn't take into consideration that the agent's policy is changing. Therefore, to account for self modification, we need to modify the second rule into ``$P(a_t\mid \mae_{< t}) = 1$ if $\pi_t(\mae_{<t})=a_t$ and zero otherwise''.  To evaluate the V and Q-functions for self-modifying agents, we need to use probabilities of complete histories calculated in this way.
	
	\begin{definition} \label{def:mod_indep}
		Let $\breve{\mae}_{1:t}$ denote the history $\mae_{1:t}$ with information about self-modification removed so that $\breve{\mae}_{1:t}=\breve{a}_1 e_1\breve{a}_2 e_2 ... \breve{a}_{t} e_t$. A function $ f \colon (\mathcal{A \times E})^* \to (anything)$ is \textit{modification-independent} if $\breve{\mae}_{1:t}=\breve{\mae}_{1:t}^{'}$ implies that $f(\mae_{1:t})=f(\mae_{1:t}^{'})$.
	\end{definition}

\paragraph{Modification-independence assumption:}In the rest of the paper, we will assume that the agent's belief and utility function as well as the correct belief are modification-independent.
